\chapter{核函数的前世今生}
\label{chap:kernel-functions}
\begin{quote}\small
\textit{本章摘要。}核函数把两个输入映射到某个特征空间后计算内积。只要这个内积能够直接由原始输入计算，许多依赖内积的线性算法就能处理原空间中的非线性关系。本章从这一计算需求出发，建立特征映射、核矩阵与正定性的联系，并比较线性、多项式、高斯径向基和拉普拉斯核的含义与局限。
\end{quote}

线性模型能否描述复杂关系，取决于输入表示。若振动幅值与故障风险呈弯曲关系，在线性模型中直接使用原始幅值可能不足；手工加入平方项可以增加表达能力，但高维输入中的交互项会迅速增多。核函数提供了另一种做法：先把输入看成处于一个更丰富的特征空间，再通过原空间中的计算得到该空间内积，而不显式构造所有特征。

\section{从特征映射到核函数}
设输入为$\bm x\in\R^d$。特征映射（feature map）$\phi$将其变为$\phi(\bm x)\in\mathcal H$，其中$\mathcal H$是带内积的特征空间。若一个算法只通过特征向量的内积使用数据，那么它所需的量是
\begin{equation}
 k(\bm x,\bm x')=\langle\phi(\bm x),\phi(\bm x')\rangle_{\mathcal H}.
 \label{eq:kernel-feature-inner-product}
\end{equation}
函数$k$称为核函数。它的输入是两个原始样本，输出是特征空间中的内积；它不必显式输出映射后的坐标。核函数的计算价值在于避免构造可能极高维甚至无限维的$\phi(\bm x)$。

先看一维映射$\phi(x)=(1,x,x^2)^{\mathsf T}$。对应内积为
\begin{equation}
 k(x,z)=1+xz+x^2z^2.
\end{equation}
例如$x=2,z=3$时，映射后分别为$(1,2,4)^{\mathsf T}$和$(1,3,9)^{\mathsf T}$，内积为$1+6+36=43$，直接代入核函数也得到$43$。线性算法若只依赖这些内积，就能在不显式建立三维特征向量的情况下完成同样计算。

\section{核函数、特征映射与特征函数不是同一个对象}
“核函数”有两个层次的含义。给定一个特征映射$\phi$时，式\eqref{eq:kernel-feature-inner-product}定义了由映射诱导的核。反过来，给定一个满足正定条件的函数$k$，可以构造至少一个特征空间和映射，使该函数成为内积。因而实际建模时常先写出$k$，不必先列出$\phi$的全部坐标。

在核方法中，特征映射$\phi(\bm x)$是一个向量或Hilbert空间元素。它把原输入变成算法真正使用的表示。特征函数$\psi_j(\bm x)$则是这个表示的第$j$个坐标，或一组标量函数。若$\phi(\bm x)=(\psi_1(\bm x),\ldots,\psi_m(\bm x))^{\mathsf T}$，则
\begin{equation}
 k(\bm x,\bm z)=\sum_{j=1}^{m}\psi_j(\bm x)\psi_j(\bm z).
 \label{eq:kernel-feature-functions}
\end{equation}
有限维时，这就是普通向量内积。无限维时，求和理解为收敛的平方可和级数。特征函数不是概率论中的特征函数$\mathbb E[e^{itX}]$，后者描述随机变量的分布，不能直接替代核方法中的特征映射。

特征映射也不唯一。若$Q$是特征空间中的等距变换，则$Q\phi(\bm x)$产生同一个核，因为$\langle Q\phi(\bm x),Q\phi(\bm z)\rangle=\langle\phi(\bm x),\phi(\bm z)\rangle$。核决定的是样本之间的内积关系，不决定某一组唯一的坐标轴。这个事实解释了为什么核PCA的特征方向可能发生符号变化或旋转，而核值和预测函数保持不变。

\section{有效核的充要条件：从有限矩阵到特征空间}
给定任意有限样本$\bm x_1,\ldots,\bm x_n$，定义Gram矩阵
\begin{equation}
 K_{ij}=k(\bm x_i,\bm x_j).
\end{equation}
核函数称为对称正定核或Mercer核的基本代数条件，是对任意$n$、任意样本和任意$\bm a\in\R^n$都有
\begin{equation}
 K=K^{\mathsf T},\qquad
 \bm a^{\mathsf T}K\bm a=\sum_{i,j}a_ia_jk(\bm x_i,\bm x_j)\geq0.
 \label{eq:kernel-psd-condition}
\end{equation}
必要性来自内积表示：
\begin{equation}
 \bm a^{\mathsf T}K\bm a
 =\left\|\sum_i a_i\phi(\bm x_i)\right\|_{\mathcal H}^2\geq0.
\end{equation}
反方向可由核诱导的函数空间构造特征映射。先取所有形式有限和$f=\sum_i a_i k(\cdot,\bm x_i)$，定义
\begin{equation}
 \left\langle\sum_i a_i k(\cdot,\bm x_i),
 \sum_j b_j k(\cdot,\bm x_j)\right\rangle
 =\sum_{i,j}a_ib_jk(\bm x_i,\bm x_j).
\end{equation}
正定性保证右侧非负，并在消去零范数函数后得到内积空间。完备化得到再生核Hilbert空间（Reproducing Kernel Hilbert Space，RKHS），其中$k(\cdot,\bm x)$本身属于该空间，并满足
\begin{equation}
 f(\bm x)=\langle f,k(\cdot,\bm x)\rangle_{\mathcal H_k}.
 \label{eq:reproducing-property}
\end{equation}
这就是“再生”一词的含义：评价函数值等价于与一个核切片做内积。它不是说核函数能生成观测数据，也不表示所有函数都具有有限范数。

半正定允许$K$有零特征值。若对所有互异样本的非零系数都有严格不等式，则称相应有限Gram矩阵严格正定。严格正定比半正定更强，但重复样本、对称样本或有限维特征映射都可能导致奇异矩阵。给定样本的数值特征值非负只能说明这一矩阵近似有效，不能证明函数对所有输入都正定。实际计算还需检查最小特征值的浮点误差。

\section{谱展开与核的构造规则}
对有限训练集，$K$可作谱分解$K=V\operatorname{diag}(\kappa_j)V^{\mathsf T}$。若$\kappa_j\geq0$，可把第$j$个训练点的特征坐标取为$\sqrt{\kappa_j}V_{ij}$，于是$K$正是这些坐标的Gram矩阵。这说明每个有限核矩阵都能找到一组等价的有限特征坐标，但这些坐标只针对当前样本，不能自动给出新输入的映射。连续输入空间上的特征映射需要由核本身或其解析展开定义。

常用的闭包规则可直接由内积或函数空间构造验证：若$k_1,k_2$有效，则$k_1+k_2$、$c k_1$（$c\geq0$）和$k_1k_2$有效；若$g$是特征空间上的线性变换，则$g(\bm x)^{\mathsf T}g(\bm z)$也是有效核。多项式核正是内积核经过乘方和加常数后的组合。用这些规则可以表达“线性相似性加局部相似性”，但组合后的尺度仍需验证。

\section{为什么要计算核}
考虑带正则化的线性预测器$f(\bm x)=\bm w^{\mathsf T}\phi(\bm x)$。当训练目标是平方误差加上$\|\bm w\|_2^2$惩罚时，最优解可写成训练特征向量的线性组合。于是预测只依赖训练样本与查询点的内积。把这些内积替换为$k(\bm x_i,\bm x_j)$，便得到核化算法。这个过程通常称为核技巧（kernel trick）。

核技巧解决的是表示与计算之间的矛盾：模型在特征空间中可以是线性的，在原始输入空间中却对应非线性函数。它不会自动发现适合任务的相似性。核的形式及其尺度参数仍表达建模选择，必须用训练数据拟合并在独立验证数据上选择。核矩阵的存储和分解通常需要$O(n^2)$内存与$O(n^3)$时间，样本很大时也会成为限制。

\section{常见核及其含义}
线性核$k(\bm x,\bm z)=\bm x^{\mathsf T}\bm z$保留原特征空间中的线性关系，计算简单，适合高维稀疏数据。它的相似性受特征尺度影响，因此单位不同的变量通常需要按训练集统计量标准化。

多项式核$k(\bm x,\bm z)=(\gamma\bm x^{\mathsf T}\bm z+c)^p$在适当参数条件下对应有限阶多项式特征。次数$p$控制交互阶数，较高次数能表达更复杂关系，也更容易对有限样本过拟合。$\gamma,c$改变各阶项的相对权重，必须与输入尺度一并考虑。

高斯径向基函数核（Gaussian radial basis function kernel，RBF核）常写为
\begin{equation}
 k(\bm x,\bm z)=\exp\left(-\frac{\|\bm x-\bm z\|_2^2}{2\ell^2}\right),\qquad \ell>0.
\end{equation}
相同输入的核值为1，距离增大时核值下降。长度尺度$\ell$控制相似性衰减的范围：小$\ell$使影响局部化，大$\ell$使较远样本仍相互关联。高斯核对应无限维特征空间，但这不表示数据量无限，也不保证拟合优良。拉普拉斯核$k(\bm x,\bm z)=\exp(-\|\bm x-\bm z\|_1/\ell)$以$L_1$距离控制衰减，对各坐标差异的累积方式不同于RBF核。

核的选择应对应可说明的相似性假设。若特征的整体欧氏距离有意义，可从RBF核开始；若领域知识指出周期或时间结构，应选择能表达该结构的核或组合核。不同核可以相加或乘：非负加权和仍为有效核，两个有效核的乘积也仍为有效核。核参数应只用训练折拟合，避免验证和测试数据的信息泄漏。

\section{核方法的边界}
核方法把相似性假设写进函数空间。若输入距离不能反映目标关系，例如未缩放的量纲支配了距离，核矩阵就会编码错误的邻近关系。核方法也通常需要保存训练样本或核矩阵，预测成本随样本数增加。核函数提供一种可控的非线性表示，不会替代数据检查、模型选择和外部验证。

后文按依赖关系展开。支持向量机先展示核如何进入分类器的对偶问题。随后同一章的PCA、PPCA与kPCA依次从原空间的几何压缩、概率生成模型和核特征空间降维展开。核岭回归、支持向量回归与核最小均方分别处理批量平方损失、容忍带损失和在线更新。高斯过程最后用核表达函数值的协方差，并与核岭回归的点预测作比较。
